\section{扩展模型与方案稳健性}

竞赛题通常给出确定性的到货量和处理效率，但实际分拣中心面对的是预测误差、设备波动和人员临时缺勤。若只报告确定性最优解，方案在数据发生小幅变化时可能需要大规模调整。本节在保持原模型主线不变的基础上，引入安全余量、场景鲁棒性和公平性目标，说明如何把本文结果扩展为具有实际部署价值的决策工具。

\subsection{预测误差下的安全容量}

令$\tilde a_{dh}=a_{dh}+\xi_{dh}$表示含误差的实际到货量，其中$\xi_{dh}$为预测误差。最直接的做法是在模型中设置小时安全容量$s_{dh}$，将容量约束改写为
\begin{equation}
 p_{dh}+s_{dh}\le \sum_s\delta_{sh}
 \left(25y_{ds}-15z_{ds,h-s}\right).
\end{equation}
安全容量不应在所有小时平均分配，而应优先配置在历史误差方差较大的高峰时段。若可估计误差标准差$\sigma_{dh}$，可以令$s_{dh}=\kappa\sigma_{dh}$，其中$\kappa$由期望服务水平确定。这样，服务承诺不再依赖单一预测值，而是对一定范围内的需求波动具有缓冲。

另一方面，安全容量会增加人数，因此需要把风险成本和用工成本放在同一目标中：
\begin{equation}
 \min\; \sum_{d,s}y_{ds}
 \lambda\sum_{d,h}s_{dh}
 \gamma\sum_{d,h}b_{dh}.
\end{equation}
第一项控制人员规模，第二项控制为不确定性预留的能力，第三项控制日内积压。$\lambda$越大，模型越倾向于通过增加早班人员获得稳定性；$\gamma$越大，模型越倾向于提前处理货物，减少库存峰值。基准模型相当于把安全容量和库存惩罚系数设为零，但仍保留硬性的清零和截止约束。

\begin{figure}[htbp]
\centering
\includegraphics[width=.88\textwidth]{figures/capacity_utilization.pdf}
\caption{基准利用率曲线为安全容量配置提供参照}
\label{fig:robust-utilization}
\end{figure}

图中利用率已经接近1的小时不适合继续承担预测上调风险；利用率较低但到货波动较大的小时也不能简单视为可削减时段。安全容量的配置应同时参考利用率和误差波动，体现“能力紧张程度”和“需求不确定程度”的乘积。

\subsection{多场景鲁棒排班}

设有$K$个需求场景，场景$k$的到货量为$a_{dh}^{(k)}$。若班次决策在场景发生前确定，可令$x_{ds}$和$y_{ds}$在所有场景之间共用，而允许处理量$p_{dh}^{(k)}$和库存$b_{dh}^{(k)}$随场景变化。于是容量约束对每个场景均成立：
\begin{equation}
 0\le p_{dh}^{(k)}
 \le \sum_s\delta_{sh}
 (25y_{ds}-15z_{ds,h-s}^{(k)}),
\quad k=1,\ldots,K.
\end{equation}
如果低效率小时属于人员行为而非需求场景，也可以让$z$在场景间共用，从而得到更保守的部署方案。场景模型的目标可以是最坏场景目标，也可以是期望目标：
\begin{equation}
 \min \; \sum_{d,s}y_{ds}
 \beta\max_k\sum_{d,h}b_{dh}^{(k)}
\left(1-\beta\right)\sum_k\pi_k
\sum_{d,h}b_{dh}^{(k)}.
\end{equation}
其中$\pi_k$为场景概率，$\beta$控制保守程度。$\beta=1$强调极端风险，$\beta=0$强调平均表现。本文基准数据可直接作为中心场景，再从历史误差或上下浮动比例构造其余场景。

\begin{table}[htbp]
\centering
\caption{确定性模型与场景模型的决策差异}
\label{tab:robust-compare}
\begin{tabular}{p{.23\textwidth}p{.26\textwidth}p{.35\textwidth}}
\toprule
比较维度 & 确定性基准 & 多场景扩展\\
\midrule
班次决策 & 针对一条到货序列优化 & 场景发生前共用一组班次\\
流量决策 & 一组处理量和库存 & 每个场景分别回代处理量和库存\\
风险控制 & 由硬约束保证基准可行 & 控制最坏场景或加权风险\\
编制解释 & 581人为基准峰值下界 & 编制随风险容忍度变化\\
\bottomrule
\end{tabular}
\end{table}

场景扩展的一个重要优点是可以把“是否增加固定编制”的问题量化。若增加一名固定员工的成本低于高需求场景下的临时工、延期处理或服务违约成本，则应选择更高的稳健程度；反之，可维持581人的基准编制并用临时人员吸收尾部需求。

\subsection{设备、技能与班次成本}

本文假设所有工人的处理效率相同。若中心存在不同技能等级或不同设备类型，可将人数变量扩展为$y_{dsg}$，其中$g$表示技能组；处理能力改写为
\begin{equation}
 p_{dh}\le\sum_{s,g}\delta_{sh}c_g y_{dsg}
 -\sum_{s,g}\delta_{sh}\ell_g z_{dsg,h-s},
\end{equation}
其中$c_g$为技能组正常效率，$\ell_g$为低效率损失。若某类货物只能由特定技能组处理，还应增加货物流分层和技能覆盖约束。这样的扩展会放大模型规模，但不改变库存守恒的基本结构。

班次成本也可以显式加入。设$c_s$为起点为$s$的单位人员成本，$o_{dh}$为小时加班或延期惩罚，则目标为
\begin{equation}
 \min\sum_{d,s}c_s y_{ds}
 +\mu\sum_{d,h}o_{dh}b_{dh}
 +\nu J_{\mathrm{fair}}.
\end{equation}
这使得模型能够在人数相同的多组方案中选择夜班更少、库存更低或员工切换更平滑的方案。对于竞赛论文，先给出人数最小的主结果，再把成本、公平性和稳定性作为二级目标进行对比，能够增强结果的完整性。

\begin{figure}[htbp]
\centering
\includegraphics[width=.88\textwidth]{figures/daily_demand_workers.pdf}
\caption{日需求与人数曲线为编制和临时工分界提供依据}
\label{fig:robust-demand}
\end{figure}

图中需求低于月度峰值的日期可以承担轮休和机动人员调配；接近峰值的日期则需要优先保障覆盖。若加入班次成本后最优起点发生变化，应同时报告成本变化和服务余量，避免只展示一个更低的目标值而忽略夜班或库存代价。

\subsection{缺勤与应急修复}

令$q_{id}=1$表示工人$i$在第$d$天临时不可用，则出勤约束可改写为
\begin{equation}
 w_{id}\le1-q_{id}.
\end{equation}
当缺勤人数较少时，可以固定其余日期的出勤决策，仅对受影响窗口重新求解；当缺勤发生在峰值日时，则需要同时检查$N=581$的可行性。若不可行，模型可以引入临时工变量$v_d$：
\begin{equation}
 \sum_iw_{id}+v_d\ge r_d,\qquad
 \min N+\tau\sum_dv_d.
\end{equation}
较大的$\tau$会优先维持固定编制，只有无法覆盖时才使用临时工；较小的$\tau$则允许以临时工换取更少的固定人员。此机制把“应急加人”从经验判断转化为带成本的可解释决策。

\begin{table}[htbp]
\centering
\caption{应急场景下的修复优先级}
\label{tab:emergency}
\begin{tabular}{p{.24\textwidth}p{.28\textwidth}p{.34\textwidth}}
\toprule
故障类型 & 首先固定的约束 & 优先采取的动作\\
\midrule
普通日少量缺勤 & 日覆盖和日末清零 & 调整机动人员与班次分配\\
峰值日缺勤 & 峰值日覆盖和16时截止 & 启用临时工并提前早班\\
连续多日缺勤 & 23工日和滚动窗口 & 重排相邻休息日\\
设备效率下降 & 小时容量和服务截止 & 更新效率参数后重新求解\\
\bottomrule
\end{tabular}
\end{table}

\subsection{扩展结果的评价原则}

扩展模型不应只比较目标值，还应比较可行率、服务余量、人员稳定性和解释成本。设场景$k$下的截止违约指示为$I_k$，则稳健性可用
\begin{equation}
 P_{\mathrm{service}}=1-\frac{1}{K}\sum_{k=1}^{K}I_k
\end{equation}
衡量；人员稳定性可用相邻滚动方案的出勤变化量衡量。理想方案是在固定编制增加有限的情况下，提高$P_{\mathrm{service}}$并减少大规模换班。本文基准方案已经完成确定性约束的全量复核，未来引入场景后可沿用同一证据注册和发布审计流程。

本节说明，本文的框架并不依赖某一组固定数字。只要输入仍然可以表示为带时间索引的到货量，处理能力可以写成可加形式，人员规则可以写成线性或整数约束，就能够继续扩展到鲁棒排班、技能排班和应急重排。扩展后的模型应保持与基准模型相同的结果链路：每个新增参数都有来源，每个新增情景都有定义，每个最终结论都有独立验证。

\clearpage
\subsection{参数校准与复用接口}

为了使框架能够持续产生不同题目的论文，模型参数不应散落在正文段落中，而应集中为可替换接口。对本题，数据接口是日期、小时和到货量；能力接口是正常效率、低效率效率和班次长度；规则接口是班次数量、服务截止、月度工日和连续工作窗口。更换数据集时只需更新接口及其来源，模型的库存守恒、覆盖约束和证据注册结构均可保留。

\begin{table}[htbp]
\centering
\caption{可复用论文框架的参数接口}
\label{tab:parameter-interface}
\begin{tabular}{p{.20\textwidth}p{.28\textwidth}p{.34\textwidth}}
\toprule
接口类别 & 本题取值或形式 & 更新时的验证要求\\
\midrule
时间接口 & 30天、每24小时 & 联合键完整且无重复\\
班次接口 & 5个连续8小时班次 & 起点不越过自然日\\
能力接口 & 正常25件/小时，低效10件/小时 & 与效率情景保持单位一致\\
服务接口 & 0--12时货物在16时前处理 & 逐日检查累计处理量\\
人员接口 & 每人23个工日，连续8日最多7日 & 逐人检查工日与窗口\\
证据接口 & JSON、图表、PDF和运行清单 & 路径在项目内且哈希可追溯\\
\bottomrule
\end{tabular}
\end{table}

参数接口的价值在于把“重新写一篇论文”变为“重新运行一条经过验证的流水线”。新题目只要补充题目陈述、数据审计、问题映射、模型注册、结果注册、图表注册和验证记录，框架就能自动生成结构化报告、编译论文并检查正文页数。若某个新题目不满足当前的连续班次假设，系统应在模型注册中明确声明限制，而不能静默套用旧模型。
